% STATUS: DRAFT
\chapter{Crossed product for black holes}\label{ch:bh_crossed}
\skippable

\begin{physmotiv}
For the large-$N$ black hole, the boundary Hamiltonian $H$ (the ADM energy) is a physical operator, not a constraint. Witten's observation is that this is precisely what turns the type III$_1$ bulk exterior algebra into a type II$_\infty$ one: the algebra of the boundary at large $N$ is generated by the bulk exterior fields \emph{and} $H$ (with $H-\ldots$ the dressing), i.e.\ it is a crossed product. Chandrasekaran, Penington and Witten computed the entropy and showed it is the generalized entropy.
\end{physmotiv}

\section{Witten's construction}

Let $\A$ be the type III$_1$ algebra of bulk fields in the exterior of the eternal black hole in the Hartle--Hawking state, with modular flow $\Ad\ee^{-\ii\beta tH_{\rm bulk}}$-type (the Schwarzschild boost). Gauge invariance under diffeomorphisms that act at the boundary means that bulk operators must be dressed to the boundary, and the boundary Hamiltonian $H$ (which has a fluctuating spectrum in the large-$N$ limit of order $N^0$ around the classical mass) is a new operator in the algebra. Writing $H=\frac{1}{\beta}K+\ldots$ with $K$ the modular Hamiltonian of the bulk fields up to a constant, the algebra is
\begin{equation}
\A\rtimes_\sigma\R\;\cong\;\{\hat a,\ H\ \text{(shifted)}\}'',
\end{equation}
with the physical Hamiltonian playing the role of $X=K+x$ in \cref{ch:crossed}. It is type II$_\infty$ (the energy spectrum is unbounded above but the trace is $\ee^{\beta H}$-weighted, \cref{ch:takesaki}). \emph{There is no maximum-entropy state}: the black hole is not in equilibrium with a heat bath in the same sense as de Sitter with a bounded-below observer; instead the entropy of a state is defined up to an additive constant and is unbounded above, as expected for a black hole that can grow.

\section{Entropy}

Chandrasekaran, Penington and Witten \cite{CPW2022} showed that the entropy of semiclassical states of the crossed product agrees, up to a state-independent constant, with the generalized entropy $\Sgen=A/4G_N+S_{\rm bulk}$; this is the first and cleanest derivation of the generalized entropy from a von Neumann entropy of a type II algebra (the formula derived in a finite-dimensional model in \cref{ch:entropy_II}).

\begin{center}
\renewcommand{\arraystretch}{1.3}
\begin{tabular}{@{}lll@{}}
\toprule
 & Black hole (large $N$) & de Sitter (observer)\\
\midrule
Clock & boundary Hamiltonian $H$, spectrum unbounded above & observer energy $q\ge0$\\
Algebra & type II$_\infty$ & type II$_1$\\
Max entropy state & none & empty dS\\
Trace weight & $\ee^{\beta H}$ & $\ee^{-\beta q}$\\
\bottomrule
\end{tabular}
\end{center}

\begin{whatwrong}
The black hole construction has a physical $H$ already in the boundary theory. In de Sitter there is none; the observer must be added. The algebraic structure is the same; the physical origin of the clock is different, and so is the sign structure that decides II$_\infty$ versus II$_1$.
\end{whatwrong}

\section*{Exercises}
\begin{exercise}
Using \cref{ch:takesaki}, show that the trace on the black hole algebra is $\tau(\hat a)=\int\dd x\,\ee^{x}\ldots$ with $x\sim\beta(H-\bar H)$, and that $\tau(\id)=\infty$ because $x$ has no upper bound.
\end{exercise}
\begin{exercise}
Compare the two-column table above: which sign flip converts $\ee^{\beta H}$ to $\ee^{-\beta q}$ in the trace weight, and what does this say about the direction in which energy can be added?
\end{exercise}
